\documentclass[10pt,a4paper]{amsart}
\usepackage[margin=19mm]{geometry}
\usepackage{amsmath,amssymb,mathtools}
\usepackage[scr=rsfs,cal=euler]{mathalfa}
\usepackage{xfrac}
\usepackage[unicode,hidelinks,bookmarks=false]{hyperref}
\newcommand{\ie}{i.e.\ }
\newcommand{\cf}{cf.\ }
\newcommand{\resp}{respectively\ }
\newcommand{\Subsection}{\subsection}
\providecommand{\nobreakdash}{\nobreak\hbox{-}}
\setlength{\emergencystretch}{2em}
\multlinegap0pt
\usepackage{bm}
\usepackage{bbm}
\usepackage{dsfont}
\usepackage{xspace}
\usepackage{cancel}
\usepackage{mathtools}
\usepackage{amsthm}
\makeatletter
\@ifundefined{theorem}{\newtheorem{theorem}{Theorem}}{}
\@ifundefined{lemma}{\newtheorem{lemma}[theorem]{Lemma}}{}
\makeatother
\def\R{\mathbb{R}}
\newcommand{\C}{\mathbb{C}}
\newcommand{\ga}{\mathfrak{a}}
\title[A truncated inner product formula in the geometry of numbers]{A truncated inner product formula in the geometry of numbers}
\author{Seokho Jin, Seungki Kim}
\date{Source: arXiv:2503.20010v3 (CC BY 4.0)}
\begin{document}
\maketitle

\noindent\emph{Source and licence.} A truncated inner product formula in the geometry of numbers. Version arXiv:2503.20010v3 (CC BY 4.0), distributed under the Creative Commons Attribution 4.0 International licence, as recorded in the arXiv licence metadata for this exact version and shown on the abstract page https://arxiv.org/abs/2503.20010v3. Full rights notice: licenses/arxiv-2503.20010-CC-BY-4.0.txt.\par\smallskip

\noindent\textit{Editorial scope.} Counted unit: the Theorem whose source label is \texttt{thm:truncate} in the article (source0.tex lines 818--820, source label \texttt{thm:truncate}), delivered together with the article's own complete original proof of it (source0.tex lines 821--901, one \texttt{proof} environment of 81 lines). No mathematics is written, repaired or re-derived: statement and proof are the article's own text line for line. Required context retained uncounted: eq:arthur\_t\_defn2 (lines 573--575); lemma:umg1 (lines 709--711); lemma:umg1' (lines 717--719); lemma:umg2 (lines 724--730); lemma:umg3 (lines 764--770). Every definition, display and label of those lines is retained unchanged. Omitted from this package and not redistributed: 1--572, 576--708, 712--716, 720--723, 731--763, 771--817, 902--2304. Nothing inside a retained excerpt is omitted or altered: every retained body is delivered exactly as the article's lines are written, so no omission receipt of any kind accompanies this package. Citations: the retained excerpts carry no citation command at all, so no bibliography entry of the article is transcribed and no reference is redistributed. Reference adaptation: none was needed and none was made; every reference inside the delivered unit resolves inside it, so no unresolved reference or citation is produced. Printed slips kept literal: no printed word, symbol, hypothesis or number of the counted statement or of its proof was altered. Layout and class: the article's own class is \texttt{amsart} with packages including latexsym, amssymb, amsmath, amsfonts, xcolor, inputenc, amsthm, amsxtra, amscd, setspace, mathrsfs, multirow, graphicx, enumerate; the wrapper is the lane's pinned \texttt{amsart} editorial wrapper at 10pt on A4, which restates the theorem-like environments the retained text needs and adds the article's own macro definitions that the retained text needs (\texttt{\textbackslash C}, \texttt{\textbackslash R}, \texttt{\textbackslash ga}), with extra wrapper packages (bm, bbm, dsfont, xspace, cancel, mathtools). The delivered numbering is the unit's own. Assets: no retained span contains a figure environment, a graphics-inclusion command, a PDF-page inclusion, a table or a TikZ picture; no image file is copied into this package, the package declares no asset dependency and no image file is redistributed with it. Licence: version arXiv:2503.20010v3 of this article is distributed under the Creative Commons Attribution 4.0 International licence, as recorded in the arXiv licence metadata for this exact version and shown on the abstract page https://arxiv.org/abs/2503.20010v3. A full rights notice is delivered at licenses/arxiv-2503.20010-CC-BY-4.0.txt. Characters: every retained excerpt is pure ASCII, so the delivered engine, XeLaTeX with its default Latin Modern text font, reports no missing character.
\begin{equation}\label{eq:arthur_t_defn2}
\Lambda^C\phi(x) = \sum_{P_1 \subseteq P_2} \sum_{\delta \in \Gamma \cap P_1 \backslash \Gamma} \chi_{P_1,P_2}^C(\delta x)\phi_{P_1,P_2}(\delta x).
\end{equation}

\begin{lemma} \label{lemma:umg1}
The support of $F^{P_1}(x, C)$ is bounded on $\ga^{P_1}$. More precisely, any $x$ in the support has $\varpi(x)$ bounded by a linear function on $\log T$ for all $\varpi \in \widehat\Delta^{P_1}$.
\end{lemma}

\begin{lemma} \label{lemma:umg1'}
Let $x = x_0 + x_1 \in \ga_{P_1}$ with $x_0 \in \ga_{P_1}^{P_2}, x_1 \in \ga_{P_2}$, such that $\sigma_{P_1}^{P_2}(x) = 1$. Then $\|x_1\| \ll \|x_0\|$, where the implied constant depends only on $P_1$ and $P_2$.
\end{lemma}

\begin{lemma} \label{lemma:umg2}
If $\phi:\Gamma \backslash G \rightarrow \C$ is $\lambda$-UMG, then 
%$\phi_{P_1,P_2}$ decays rapidly along $\ga_{P_1}^{P_2}$, i.e. 
\begin{equation*}
|\phi_{P_1,P_2}(g)| \ll a(g)^{\lambda-\mu}
\end{equation*} for any linear combination $\mu$ of the elements in $\Delta^{P_2} - \Delta^{P_1}$. The implied constant depends on $\phi$ and $\mu$ only.
\end{lemma}

\begin{lemma} \label{lemma:umg3}
Suppose $g \in \mathrm{supp}\,\chi_{P_1,P_2}^C$. Then there exists a constant $p_0 \in \R$ and a linear combination $\kappa$ of $\Delta^{P_2} - \Delta^{P_1}$, both independent of $g$, such that
\begin{equation*}
\sum_{\delta \in \Gamma \cap P_1 \backslash \Gamma} \chi_{P_1,P_2}^C(\delta g) \ll T^{p_0} a(g)^\kappa.
\end{equation*}
The implied constant depends only on $G$.
\end{lemma}

\begin{theorem} \label{thm:truncate}
For $\lambda$-UMG $\phi$ and any $p > 0$, there exists a constant $c=c(p,\phi)>0$ such that $|\Lambda^C\phi(g)| < cT^{-p}$ outside the support of $\chi_{G,G}^C$.
\end{theorem}

\begin{proof}
In the light of \eqref{eq:arthur_t_defn2} and Lemma \ref{lemma:umg3}, it suffices to prove, for each $P_1 \subseteq P_2$, $P_1 \neq G$, and for $g$ contained in the support of $\chi_{P_1,P_2}^C$, that
\begin{equation*}
|\phi_{P_1,P_2}(g)| \ll T^{-p}a(g)^{-\kappa}
\end{equation*}
for any choice of $p>0$ and a certain $\kappa \in \mathrm{span}(\Delta^{P_2}-\Delta^{P_1})$.

On the other hand, we already know from Lemma \ref{lemma:umg2} that
\begin{equation*}
|\phi_{P_1,P_2}(g)| \ll a(g)^{\lambda-\mu}
\end{equation*}
holds for some $\lambda \in \ga^*$ and any $\mu \in \mathrm{span}(\Delta^{P_2} - \Delta^{P_1})$.
Write
\begin{align*}
H(g) &= H_0^2 + H_2, \ H_0^2 = H_0^1 + H_1^2, \ & H_i^j \in \ga_{P_i}^{P_j}, \\
\lambda-\mu &= \nu_0^2 + \nu_2, \ \nu_0^2 = \nu_0^1 + \nu_1^2, \ & \nu_i^j \in (\ga_{P_i}^{P_j})^*.
\end{align*}
Our goal is then to show that, for an appropriate choice of $\mu$,
\begin{equation*}
(\lambda-\mu)(H(g)) = \nu_0^1(H_0^1) + \nu_1^2(H_1^2) + \nu_2(H_2)
\end{equation*}
is bounded from above by $-p\log T$. Note that there is no loss of generality in assuming $\kappa=0$, by replacing $\mu$ with $\mu - \kappa$.

%The proof is short but a bit technical, so let us outline our idea first. Consider
%\begin{equation*}
%(\lambda-\mu)(H(g)) = \nu_0^1(H_0^1) + \nu_1^2(H_1^2) + \nu_2(H_2).
%\end{equation*}
%Roughly speaking, the support of $\chi_{P_1,P_2}^C$ is shaped so that the coordinates of $H_0^1$ are $O(\log T)$ (Lemma \ref{lemma:umg1}), those of $H_1^2$ are $\Omega(\log T)$, and $\|H_2\| \ll \|H_1^2\|$ (Lemma \ref{lemma:umg1'}). On the other hand, incrementing $m_\alpha$ leaves $\nu_2$ fixed, but ``decreases'' $\nu_1^2$ to a large negative; its effect on $\nu_0^1$ is \emph{a priori} unclear and needs to be investigated. We would like to exploit these properties to show that the above expression is dominated by $\nu_1^2(H_1^2)$, which is in turn bounded above by an arbitrarily large negative multiple of $\log T$. Then an application of Lemma \ref{lemma:umg2} will complete the proof.


Fix an $\alpha = \alpha_i \in \Delta^{P_2}-\Delta^{P_1}$, and suppose the Levi block of $P_1$ containing the $i$-th row has blocksize $n_1$, and the block containing the $(i+1)$-st row has blocksize $n_2$. We claim that
\begin{equation}\label{eq:bowwow}
\alpha(H_0^2) > \frac{1}{2}\log T + \frac{1}{n_1+n_2}\alpha(H_1^2).
\end{equation}
This follows from the additional claims that
\begin{equation} \label{eq:bowwow2}
\alpha(H_0^1) > \left(1-\frac{n_1+n_2}{2}\right)\log T, \ \alpha(H_1^2) > \left(\frac{n_1+n_2}{2}\right)\log T,
\end{equation}
by adding the former and $(1-(n_1+n_2)^{-1})$ times the latter.
To see the former inequality, observe that $F^{P_1}(g,C) = 1$ implies
\begin{equation*}
\varpi(H_0^1) < \varpi(C) = \log T \cdot \varpi((\rho^\vee)^{P_1})
\end{equation*}
for all $\varpi \in \widehat\Delta_0^{P_1}$. $(\rho^\vee)^{P_1}$ here can be computed explicitly: if $P_1$ consists of the Levi blocks of size $k_1,\ldots,k_r$, then
\begin{equation*}
(\rho^\vee)^{P_1} = \left(\frac{k_1-1}{2}, \ldots, \frac{1-k_1}{2}; \ldots ; \frac{k_r-1}{2}, \ldots, \frac{1-k_r}{2} \right).
\end{equation*}
This implies that the $i$-th entry of $H_0^1$ is greater than $-1/2 \cdot (n_1-1)\log T$, and the $(i+1)$-st entry of $H_0^1$ is less than $1/2 \cdot (n_2-1)\log T$, hence the desired inequality. Similarly, the latter inequality of \eqref{eq:bowwow2} follows from our assumption $\sigma_{P_1}^{P_2}(H_{P_1}(g) - C) = 1$, which implies that $\tau_{P_1}^{P_2}(H_1^2 - C) = 1$ and thus
\begin{equation*}
\alpha(H_1^2) > \alpha(C_{P_1}^{P_2}) = \log T \cdot \alpha((\rho^\vee)_{P_1}^{P_2}).
\end{equation*}
It can be computed that $\alpha((\rho^\vee)_{P_1}^{P_2}) = 1/2 \cdot (n_1 + n_2)\log T$, proving \eqref{eq:bowwow2} in full.

To complete the proof of the theorem, let us start by writing
\begin{align*}
\lambda_0^2 := \lambda_{P_0}^{P_2} = \lambda' + \sum_{\alpha \in \Delta^{P_2} - \Delta^{P_1}} l_\alpha \alpha
\end{align*}
for some $\lambda' \in \ga^{P_1}$ and $l_\alpha$'s --- in fact, they are uniquely determined --- and
\begin{align*}
\mu = \sum_{\alpha \in \Delta^{P_2}-\Delta^{P_1}} m_\alpha \alpha,
\end{align*}
with $m_\alpha$'s to be determined soon. Take $m_\alpha$'s large enough so that $m'_\alpha := m_\alpha - l_\alpha \geq 0$.
Then by \eqref{eq:bowwow}, there exists a constant $c(P_1,P_2)>0$ such that
\begin{align*}
\nu_0^2(H_0^2) &= \lambda_0^2(H_0^2) - \mu(H_0^2) \\
&< \lambda'(H_0^1) -\frac{1}{2}\sum_{\alpha \in \Delta^{P_2} - \Delta^{P_1}} m'_\alpha\log T - c(P_1,P_2)\sum_{\alpha \in \Delta^{P_2} - \Delta^{P_1}} m'_\alpha\alpha(H_1^2) \\
&< O_\lambda(\log T) -\frac{1}{2}\sum_{\alpha \in \Delta^{P_2} - \Delta^{P_1}} m'_\alpha\log T - c(P_1,P_2)\sum_{\alpha \in \Delta^{P_2} - \Delta^{P_1}} m'_\alpha\alpha(H_1^2).
\end{align*}
The last inequality here is a consequence of Lemma \ref{lemma:umg1}.
It remains to show that, by increasing $m_\alpha$'s if necessary, 
\begin{equation*}
c(P_1,P_2)\sum_{\alpha \in \Delta^{P_2} - \Delta^{P_1}} m'_\alpha\alpha(H_1^2) + O(\log T) > \nu_2(H_2).
\end{equation*}
Since $\nu_2$ is a linear functional, we have $|\nu_2(H_2)| \ll \|H_2\|$. Also, by Lemma \ref{lemma:umg1'} we have $\|H_2-C_2\| \ll \|H_1^2-C_1^2\|$, which implies that $\|H_2\| \ll \|H_1^2\|+O(\log T)$ by the triangle inequality.
On the other hand, $\Delta^{P_2} - \Delta^{P_1}$ induces a coordinate system on $\ga_{P_1}^{P_2}$, so $\|H_1^2\| \ll \sum_{\alpha \in \Delta^{P_2} - \Delta^{P_1}} |\alpha(H_1^2)|$; but since all $\alpha(H_1^2) > 0$ by \eqref{eq:bowwow2}, we can take the absolute value sign off the latter. Summarizing, we obtain
\begin{equation*}
|\nu_2(H_2)| \ll \sum_{\alpha \in \Delta^{P_2} - \Delta^{P_1}} \alpha(H_1^2) + O(\log T),
\end{equation*}
as desired.
This completes the proof.
\end{proof}

\end{document}
